\chapter{The Maximum Principle}

\section*{3.1 \quad Statement of the maximum principle}

\begin{theorem}
\label{topping-maxpr-weak-maximum-principle}
\dcref{3.1.1}
\group{maximum-principle}
\source{topping:page-0045,topping:page-0046}
Let $g(t)$ be a smooth family of metrics and $X(t)$ a smooth family of vector
fields on a closed manifold $\M$, for $t \in [0,T]$ with $0 < T < \infty$. Let
$F:\RR \times [0,T] \to \RR$ be smooth, and let
$u \in C^\infty(\M \times [0,T], \RR)$ satisfy
\[
\frac{\partial u}{\partial t}
\le \Delta_{g(t)} u + \langle X(t), \nabla u \rangle + F(u,t).
\]
Suppose $\phi:[0,T]\to \RR$ solves
\[
\begin{cases}
\dfrac{d\phi}{dt} = F(\phi(t),t),\\
\phi(0)=\alpha \in \RR.
\end{cases}
\]
If $u(\cdot,0)\le \alpha$, then $u(\cdot,t)\le \phi(t)$ for all $t\in[0,T]$.
\end{theorem}

\begin{proof}
For $\e>0$, consider the perturbed ODE
\[
\begin{cases}
\dfrac{d\phi_\e}{dt} = F(\phi_\e(t),t)+\e,\\
\phi_\e(0)=\alpha+\e.
\end{cases}
\]
Basic ODE theory gives a solution $\phi_\e$ on $[0,T]$ for all sufficiently
small $\e>0$, and $\phi_\e \to \phi$ uniformly as $\e \downarrow 0$.

It therefore suffices to prove $u(\cdot,t) < \phi_\e(t)$ for every
$t\in[0,T]$. If not, choose the earliest $t_0\in(0,T]$ where this fails and a
point $x\in \M$ with $u(x,t_0)=\phi_\e(t_0)$. At $(x,t_0)$ we have
\[
\frac{\partial u}{\partial t}(x,t_0)-\phi_\e'(t_0)\ge 0,
\]
while $x$ is a spatial maximum of $u(\cdot,t_0)$, so
\[
\Delta u(x,t_0)\le 0,\qquad \nabla u(x,t_0)=0.
\]
Combining these with the differential inequality for $u$ and the ODE for
$\phi_\e$ gives
\[
0 \ge \left[\frac{\partial u}{\partial t}-\Delta u-\langle X,\nabla u\rangle
-F(u,\cdot)\right](x,t_0)
\ge \phi_\e'(t_0)-F(\phi_\e(t_0),t_0)=\e>0,
\]
a contradiction.
\end{proof}

\begin{corollary}
\label{topping-maxpr-weak-minimum-principle}
\dcref{3.1.2}
\group{maximum-principle}
\source{topping:page-0045}
Theorem~3.1.1 also holds with all inequalities reversed.
\end{corollary}

\begin{remark}
\label{topping-maxpr-strong-maximum-principle}
\dcref{3.1.3}
\group{maximum-principle}
\source{topping:page-0046}
The strong maximum principle for scalars gives the strict inequality
$u(\cdot,t)<\phi(t)$ for all $t\in(0,T]$, unless $u(x,t)=\phi(t)$ for all
$x\in \M$ and $t\in[0,T]$.
\end{remark}

\section*{3.2\ \ \ Basic control on the evolution of curvature}

By applying the maximum principle to curvature evolution equations, one obtains
preliminary control on the evolution of $R$ and $\Rm$. Chapter~9 refines these
estimates.

\begin{theorem}
\label{topping-maxpr-scalar-lower-bound}
\dcref{3.2.1}
\group{maximum-principle}
\source{topping:page-0047}
\uses{topping-maxpr-weak-minimum-principle}
Suppose $g(t)$ is a Ricci flow on a closed manifold $\M$, for $t\in[0,T]$.
If $R\ge \alpha\in\RR$ at time $t=0$, then for all $t\in[0,T]$,
\[
R \ge \frac{\alpha}{1-\left(\frac{2\alpha}{n}\right)t}.
\]
\end{theorem}

\begin{proof}
Apply the weak minimum principle to the scalar curvature evolution equation
(2.5.6) with $u\equiv R$, $X\equiv 0$, and
$F(r,t)\equiv \frac{2}{n}r^2$. The corresponding ODE solution is
\[
\phi(t)=\frac{\alpha}{1-\left(\frac{2\alpha}{n}\right)t}.
\]
\end{proof}

\begin{corollary}
\label{topping-maxpr-scalar-preserved}
\dcref{3.2.2}
\group{maximum-principle}
\source{topping:page-0047}
If $g(t)$ is a Ricci flow on a closed manifold $\M$, for $t\in[0,T]$, and
$R\ge \alpha\in\RR$ at $t=0$, then $R\ge \alpha$ for all $t\in[0,T]$.
\end{corollary}

\begin{corollary}
\label{topping-maxpr-positive-scalar-preserved}
\dcref{3.2.3}
\group{maximum-principle}
\source{topping:page-0047}
Positive, or weakly positive, scalar curvature is preserved under Ricci flow on
a closed manifold.
\end{corollary}

\begin{corollary}
\label{topping-maxpr-finite-time-blowup}
\dcref{3.2.4}
\group{maximum-principle}
\source{topping:page-0047}
Suppose $g(t)$ is a Ricci flow on a closed manifold $\M$, for $t\in[0,T)$.
If $R\ge \alpha>0$ at time $t=0$, then $T\le \frac{n}{2\alpha}$.
\end{corollary}

\begin{corollary}
\label{topping-maxpr-negative-scalar-lower-bound}
\dcref{3.2.5}
\group{maximum-principle}
\source{topping:page-0047}
Suppose $g(t)$ is a Ricci flow on a closed manifold $\M$, for $t\in(0,T]$.
Then
\[
R \ge -\frac{n}{2t}
\]
for all $t\in(0,T]$.
\end{corollary}

\begin{corollary}
\label{topping-maxpr-volume-decreases}
\dcref{3.2.6}
\group{maximum-principle}
\source{topping:page-0047}
If $g(t)$ is a Ricci flow on a closed manifold $\M$, for $t\in[0,T]$, and
$R\ge 0$ at $t=0$, then the volume $V(t)$ is weakly decreasing.
\end{corollary}

\begin{corollary}
\label{topping-maxpr-volume-ratio}
\dcref{3.2.7}
\group{maximum-principle}
\source{topping:page-0047,topping:page-0048}
\uses{topping-maxpr-scalar-lower-bound}
Suppose $g(t)$ is a Ricci flow on a closed manifold $\M$, for $t\in[0,T]$.
If $\alpha:=\inf R<0$ at time $t=0$, then
\[
\frac{V(t)}{\left(1+\frac{2(-\alpha)}{n}t\right)^{\frac n2}}
\]
is weakly decreasing, and in particular
\[
V(t)\le V(0)\left(1+\frac{2(-\alpha)}{n}t\right)^{\frac n2}.
\]
\end{corollary}

\begin{proof}
Differentiate the logarithm of the displayed ratio:
\[
\frac{\dd}{\dd t}\ln\!\left[\frac{V(t)}{\left(1+\frac{2(-\alpha)}{n}t\right)^{\frac n2}}\right]
= -\frac{1}{V}\int R\,\dd V+\frac{\alpha}{1+\frac{2(-\alpha)}{n}t}.
\]
This is bounded above by
\[
-\inf_{\M}R(t)+\frac{\alpha}{1-\frac{2\alpha}{n}t}\le 0
\]
by Theorem~3.2.1.
\end{proof}

\begin{remark}
\label{topping-maxpr-perelman-volume}
\dcref{3.2.8}
\group{maximum-principle}
\source{topping:page-0048}
If the Ricci flow exists for all $t\in[0,\infty)$, then
\[
\overline V:=\lim_{t\to\infty}\frac{V(t)}{\left(1+\frac{2(-\alpha)}{n}t\right)^{\frac n2}}
\]
exists. A rough principle in Perelman's work is that, locally, $\overline V>0$
suggests the manifold is becoming hyperbolic, while $\overline V=0$ suggests a
graph-manifold-like limit.
\end{remark}

\begin{remark}
\label{topping-maxpr-preserved-curvature-conditions}
\dcref{3.2.9}
\group{maximum-principle}
\source{topping:page-0048}
Scalar curvature is only one curvature quantity whose positivity is preserved
under Ricci flow on closed manifolds. Other preserved conditions include
positive curvature operator in any dimension, positive isotropic curvature in
four dimensions, and positive Ricci curvature in dimension three or less.
Negative curvatures are not preserved in general.
\end{remark}

\begin{proposition}
\label{topping-maxpr-curvature-norm-evolution}
\dcref{3.2.10}
\group{maximum-principle}
\source{topping:page-0048,topping:page-0049}
Under Ricci flow,
\[
\frac{\partial}{\partial t}|\Rm|^2\le \Delta |\Rm|^2-2|\nabla \Rm|^2+C|\Rm|^3,
\]
where $C=C(n)$ and $\Delta$ is the Laplace-Beltrami operator.
\end{proposition}

\begin{proof}
Since $\partial_t g^{ij}=-h^{ij}=2R^{ij}$, differentiation of
\[
|\Rm|^2=g^{ij}g^{kl}g^{ab}g^{cd}R_{ikac}R_{jlbd}
\]
gives
\[
\frac{\partial}{\partial t}|\Rm|^2
\le C|\Rm|^3+2\langle \Rm,\Delta\Rm+\Rm*\Rm\rangle
\le 2\langle \Rm,\Delta\Rm\rangle+C|\Rm|^3.
\]
Since $d|\Rm|^2=2\langle \Rm,\nabla\Rm\rangle$, we also have
\[
\Delta|\Rm|^2 = 2|\nabla\Rm|^2 + 2\langle \Rm,\Delta\Rm\rangle.
\]
Weakening the inequality to
\[
\frac{\partial}{\partial t}|\Rm|^2 \le \Delta|\Rm|^2+C|\Rm|^3,
\]
completes the proposition.
\end{proof}

\begin{theorem}
\label{topping-maxpr-curvature-norm-bound}
\dcref{3.2.11}
\group{maximum-principle}
\source{topping:page-0049}
Suppose $g(t)$ is a Ricci flow on a closed manifold $M$, for $t\in[0,T]$, and
that at $t=0$ we have $|\Rm|\le M$. Then for all $t\in(0,T]$,
\[
|\Rm|\le \frac{M}{1-\frac12 CMt}.
\]
\end{theorem}

\begin{proof}
Apply the weak maximum principle to the weakened inequality from
Proposition~3.2.10 with $u=|\Rm|^2$, $X\equiv 0$, $F(r,t)=Cr^{3/2}$,
$\alpha=M^2$, and
\[
\phi(t)=\frac{1}{\bigl(M^{-1}-\frac12 Ct\bigr)^2}.
\]
This ODE solution gives $|\Rm|^2\le \phi(t)$, hence
\[
|\Rm|\le \frac{M}{1-\frac12 CMt}.
\]
\end{proof}

\begin{remark}
\label{topping-maxpr-curvature-doubling}
\dcref{3.2.12}
\group{maximum-principle}
\source{topping:page-0050}
If $|\Rm|\le 1$ at $t=0$, then the waiting time before the maximum of $|\Rm|$
doubles is bounded below by a positive constant depending only on the dimension
$n$.
\end{remark}

\section*{3.3\ \ \ Global curvature derivative estimates}

We have seen that the curvature tensor satisfies a heat-type equation. The
goal now is to control higher derivatives in terms of lower derivatives across
the whole manifold.

\begin{theorem}
\label{topping-maxpr-first-derivative-estimate}
\dcref{3.3.1}
\group{maximum-principle}
\source{topping:page-0050,topping:page-0051}
Suppose $M>0$ and $g(t)$ is a Ricci flow on a closed manifold $\M^n$ for
$t\in[0,\frac1M]$. For all $k\in\mathbb N$ there exists $C=C(n,k)$ such that
if $|\Rm|\le M$ throughout $\M\times[0,\frac1M]$, then for all
$t\in[0,\frac1M]$,
\[
|\nabla^k \Rm|\le \frac{CM}{t^{k/2}}.
\]
\end{theorem}

\begin{proof}
Only the case $k=1$ is proved here. By the curvature evolution equation and
the commutation formulae,
\[
\frac{\partial}{\partial t}\Rm = \Delta\Rm+\Rm*\Rm,\qquad
\nabla(\Delta\Rm)=\Delta(\nabla\Rm)+\Rm*\nabla\Rm,
\]
and
\[
\nabla\frac{\partial}{\partial t}\Rm
=\frac{\partial}{\partial t}\nabla\Rm+\Rm*\nabla\Rm.
\]
Therefore
\[
\frac{\partial}{\partial t}\nabla\Rm
= \Delta(\nabla\Rm)+\Rm*\nabla\Rm,
\]
and hence
\[
\frac{\partial}{\partial t}|\nabla\Rm|^2
\le \Delta|\nabla\Rm|^2-2|\nabla^2\Rm|^2
+C|\Rm|\,|\nabla\Rm|^2.
\]
Set $u=t|\nabla\Rm|^2+\alpha|\Rm|^2$. Using the estimate above together with
Proposition~3.2.10, and choosing $\alpha$ sufficiently large, one gets
\[
\frac{\partial u}{\partial t}\le \Delta u + C(n)M^3.
\]
Since $u(\cdot,0)\le \alpha M^2$, the weak maximum principle with
\[
\phi(t)=\alpha M^2+CtM^3
\]
implies $u(\cdot,t)\le CM^2$, so
\[
t|\nabla\Rm|^2\le CM^2
\]
and therefore
\[
|\nabla\Rm|\le \frac{CM}{\sqrt t}.
\]
\end{proof}

\begin{corollary}
\label{topping-maxpr-higher-derivatives}
\dcref{3.3.2}
\group{maximum-principle}
\source{topping:page-0052}
Suppose $M>0$ and $g(t)$ is a Ricci flow on a closed manifold $\M^n$ for
$t\in[0,\frac1M]$. For all $j,k\in\{0\}\cup\mathbb N$, there exists
$C=C(j,k,n)$ such that if $|\Rm|\le M$ throughout
$\M\times[0,\frac1M]$, then
\[
\left|\frac{\partial^j}{\partial t^j}\nabla^k \Rm\right|
\le \frac{CM}{t^{j+\frac{k}{2}}}.
\]
\end{corollary}

\begin{proof}
By rescaling, it suffices to prove the claim at $t=1$, where the assumption
gives $M\le 1$. The cases $(j,k)=(1,0)$ and $(1,1)$ follow from the evolution
equations for $\Rm$ and $\nabla\Rm$ together with Theorem~3.3.1.

For $j=1$ and $k>1$, an induction using
\[
\frac{\partial}{\partial t}\nabla^k\Rm
= \nabla\frac{\partial}{\partial t}\nabla^{k-1}\Rm+\nabla^{k-1}\Rm*\nabla\Rm
\]
shows that $\partial_t\nabla^k\Rm$ is a $*$-composition of terms involving
$\Rm$ and its covariant derivatives, each bounded at $t=1$ by the previously
handled cases. Repeating the same idea yields the higher $j$ estimates.
\end{proof}

\begin{theorem}
\label{topping-maxpr-shi-estimate}
\dcref{3.3.3}
\group{maximum-principle}
\source{topping:page-0053}
Suppose that $g(t)$ is a Ricci flow defined on a boundaryless manifold $U$,
for $t\in[0,T]$, with no requirement that $g(t)$ is complete. Suppose further
that $|\Rm|\le M$ on $U\times[0,T]$, and that $p\in U$ and $r>0$ satisfy
$B_{g(0)}(p,r)\subset U$. Then
\[
|\nabla\Rm(p,T)|^2\le C(n)M^2\left(\frac1{r^2}+\frac1T+M\right).
\]
\end{theorem}
